\anexo{Entrevistas}
\label{app:entrevistas}

A continuación se transcriben las tres entrevistas semiestructuradas realizadas como parte del user research (véase la sección~\ref{sec:analisis_entrevistas}), a un jugador promedio, un creador de contenido de videojuegos y un estudiante de desarrollo de videojuegos.

\section*{Persona 1 (Jugador Promedio)}
\addcontentsline{toc}{section}{Persona 1 (Jugador Promedio)}

\begin{description}[leftmargin=0cm, style=nextline]
	\item[\textbf{1. ¿Cómo elegís actualmente qué videojuego jugar y qué dificultades encontrás al tomar esa decisión?}]
	Mi principal fuente para elegir un nuevo juego son los gameplays que veo a través de redes sociales, donde puedo ver de forma sencilla cómo son las mecánicas y la dinámica del juego para saber si me puede divertir, especialmente si es para jugar con amigos. Es muy difícil que hoy en día decida jugar un juego solamente por encontrarlo navegando en una tienda.

	Una de las principales dificultades es que hay una cantidad enorme de juegos disponibles y no tengo interés en investigar uno por uno. Por eso, normalmente necesito algún estímulo previo, como que me aparezca un gameplay, que un amigo me lo recomiende o que vea que varias personas de mi entorno lo están jugando. También influye mucho que mis amigos estén dispuestos a jugarlo, porque puedo encontrar un juego interesante, pero si está pensado principalmente para jugar en grupo y ninguno de mis amigos lo tiene o quiere comprarlo, probablemente termine descartándolo.

	\item[\textbf{2. ¿Qué información necesitás para saber si un videojuego puede gustarte antes de empezarlo o comprarlo?}]
	Primero necesito saber si se puede jugar con amigos o si está pensado principalmente para jugar solo, porque eso es determinante en mi decisión. Después necesito ver algún video o gameplay que me permita entender cómo se juega realmente, más allá de un tráiler o de la descripción que aparece en la tienda.

	También me sirve mucho saber si se parece a algún otro juego que ya haya jugado. Si puedo relacionarlo con juegos que sé que me gustan o no me gustan, puedo proyectar mucho mejor si voy a disfrutarlo. Por eso, para mí es más útil que me digan \enquote{se parece a este juego por estas razones} que simplemente ver el género al que pertenece o una puntuación. También considero importante saber qué tan necesario es jugarlo con otras personas, cuánto tiempo o dedicación requiere y, obviamente, el precio, especialmente si tengo dudas sobre cuánto lo voy a terminar jugando.

	\item[\textbf{3. ¿Qué tanta importancia les das a las recomendaciones de videojuegos que recibís y qué hace que decidas seguirlas o ignorarlas?}]
	Son fundamentales. Hoy casi no empiezo juegos nuevos por descubrimiento completamente propio; la mayoría de las cosas nuevas que juego surgen a partir de recomendaciones, ya sea de amigos o de gameplays y contenido que veo en redes sociales.

	Normalmente les doy mucha más importancia a las recomendaciones de mis amigos más cercanos, porque conocemos qué tipos de juegos nos gustan y, además, muchas veces la recomendación implica que vamos a poder jugarlo juntos. También tengo en cuenta las recomendaciones de creadores de contenido cuyos gustos ya conozco. Si sigo a alguien hace tiempo y sé que suele disfrutar juegos similares a los que me gustan a mí, su recomendación tiene mucho más peso que la de alguien que no conozco.

	En cambio, suelo ignorar recomendaciones demasiado genéricas o que siento que están motivadas simplemente porque un juego es nuevo, popular o está siendo promocionado. Para mí es importante entender por qué me están recomendando específicamente ese juego.

	\item[\textbf{4. Cuando una plataforma te recomienda un videojuego, ¿qué información debería mostrarte para que entiendas por qué fue elegido para vos?}]
	Principalmente me gustaría saber si mis amigos ya tienen o juegan ese juego, a qué otros juegos que yo haya jugado se parece y si está pensado para jugar en multijugador o solo. También sería útil que la plataforma relacione la recomendación con mi propio comportamiento, por ejemplo: \enquote{te recomendamos este juego porque jugaste muchas horas a X}, \enquote{porque tenés X en tu lista de deseados} o \enquote{porque varios de tus amigos están jugando}.

	No quiero que me recomiende un juego simplemente porque sea un lanzamiento nuevo, sea popular en ese momento o porque me lo quieran vender. Preferiría que mi lista de recomendaciones se construya realmente a partir de mis gustos: los juegos que jugué y cuánto los jugué, mi lista de deseados, lo que juegan mis amigos y las recomendaciones que ellos hayan hecho. Para mí, mientras más claro sea el motivo de la recomendación, más confianza tendría en ella.

	\item[\textbf{5. Si pudieras diseñar una herramienta ideal para ayudarte a elegir videojuegos, ¿qué tendría que hacer para que realmente la usaras?}]
	La herramienta ideal tendría que aprender de lo que realmente juego y no solamente de los juegos que compro o miro. Me gustaría que pudiera analizar mi historial y detectar qué características tienen en común los juegos a los que más horas les dedico, y utilizar esa información para recomendarme otros.

	También tendría que incorporar el componente social. Para mí sería muy útil poder entrar y ver directamente qué juegos podrían gustarme y que, al mismo tiempo, ya tengan algunos de mis amigos. Incluso podría diferenciar entre recomendaciones para jugar solo y recomendaciones para jugar con determinados amigos, teniendo en cuenta los juegos que tenemos en común y nuestros gustos.

	Además, cada recomendación debería estar acompañada de una explicación concreta de por qué cree que ese juego me puede gustar, mostrando similitudes con juegos que ya jugué, cuáles de mis amigos lo tienen, qué tipo de experiencia ofrece y algún gameplay corto que me permita entender rápidamente sus mecánicas.

	Para que realmente la usara tendría que ahorrarme el trabajo de buscar toda esa información por separado. La utilidad estaría en poder entrar a un solo lugar y encontrar recomendaciones personalizadas, entender rápidamente por qué me las está haciendo y tener suficiente información para decidir si vale la pena probar o comprar un juego.
\end{description}

\section*{Persona 2 (Creador de contenido de Videojuegos)}
\addcontentsline{toc}{section}{Persona 2 (Creador de contenido de Videojuegos)}

\begin{description}[leftmargin=0cm, style=nextline]
	\item[\textbf{1. ¿Cómo elegís actualmente qué videojuego jugar y qué dificultades encontrás al tomar esa decisión?}]
	A la hora de elegir un videojuego, suelo guiarme bastante por las recomendaciones de Steam y por la valoración general de los usuarios. Me llaman especialmente la atención aquellos títulos que cuentan con reseñas \enquote{extremadamente positivas}, porque siento que existe un consenso bastante fuerte sobre su calidad. También descubro muchos juegos a través de creadores de contenido que sigo y cuyos gustos conozco, como Alexelcapo, BaityBait o Knekro. Verlos jugar me permite conocer propuestas que quizá no encontraría buscando por mi cuenta.

	\item[\textbf{2. ¿Qué información necesitás para saber si un videojuego puede gustarte antes de empezarlo o comprarlo?}]
	Me interesa conocer la duración aproximada antes de comprar un juego, ya que me molesta invertir dinero en títulos que se terminan demasiado rápido. Sin embargo, tampoco busco solamente juegos largos, porque algunos pueden volverse monótonos o extenderse de manera innecesaria. Además de la duración, necesito ver algún gameplay para entender cómo se juega realmente y evaluar si sus mecánicas podrían mantenerme enganchado.

	\item[\textbf{3. ¿Qué tanta importancia les das a las recomendaciones de videojuegos que recibís y qué hace que decidas seguirlas o ignorarlas?}]
	Le doy mucha importancia a las recomendaciones de otras personas, especialmente cuando provienen de alguien que considero que tiene conocimientos o experiencia en videojuegos. Confío más en la opinión de personas que jugaron una gran variedad de títulos y pueden comparar diferentes géneros, mecánicas y experiencias. De todas formas, también tengo en cuenta si sus gustos son similares a los míos antes de seguir una recomendación.

	\item[\textbf{4. Cuando una plataforma te recomienda un videojuego, ¿qué información debería mostrarte para que entiendas por qué fue elegido para vos?}]
	Generalmente utilizo la sección de Steam que muestra juegos \enquote{recomendados para vos} según tus gustos, y considero que suele acertar. Las etiquetas de género también me resultan útiles porque permiten reconocer rápidamente las características principales de un juego. No me gusta tener que investigar durante demasiado tiempo para saber si un título puede interesarme. Para mí, lo ideal sería que todos los juegos ofrecieran una demo, ya que probarlos directamente es la mejor manera de resolver cualquier duda antes de comprarlos. Lamentablemente, no todos cuentan con esa posibilidad.

	\item[\textbf{5. Si pudieras diseñar una herramienta ideal para ayudarte a elegir videojuegos, ¿qué tendría que hacer para que realmente la usaras?}]
	Una plataforma de recomendación debería mostrar gameplays o fragmentos representativos de cada juego, informar su duración aproximada e incorporar reseñas de otros jugadores. Me gustaría poder conocer rápidamente sus principales aspectos positivos y negativos, sin tener que buscar información en distintas páginas. También sería útil que explicara por qué recomienda cada título y qué elementos comparte con otros juegos que ya disfruté.
\end{description}

\section*{Persona 3 (Estudiante de Desarrollo de Videojuegos)}
\addcontentsline{toc}{section}{Persona 3 (Estudiante de Desarrollo de Videojuegos)}

\begin{description}[leftmargin=0cm, style=nextline]
	\item[\textbf{1. ¿Cómo elegís actualmente qué videojuego jugar y qué dificultades encontrás al tomar esa decisión?}]
	Actualmente juego mayormente videojuegos con amigos, por lo que no es una decisión que tome yo solo. Pero el criterio principal diría que es el mismo: que el videojuego sea acorde a lo que siento, o tengo ganas de, en ese momento. No voy a jugar al mismo videojuego si vuelvo a mi casa cansado de la facultad y solo quiero relajarme y despejar la cabeza, que si tengo energía y ganas de jugar a algo más competitivo o desafiante.

	Por ende, la dificultad más grande que encuentro a la hora de decidir qué jugar es encontrar un videojuego que cumpla exactamente con lo que yo busco de él en ese momento.

	\item[\textbf{2. ¿Qué información necesitás para saber si un videojuego puede gustarte antes de empezarlo o comprarlo?}]
	Antes de empezar a jugar a un videojuego me guío principalmente por el tráiler, o en lo posible por un tramo de gameplay, por las reseñas, y, en algunos casos, por el desarrollador.

	En mi opinión, la mejor manera de juzgar si a uno le atrae un videojuego antes de jugarlo es viendo gameplay. Si a mí no me parece atractivo el videojuego luego de ver su jugabilidad, es muy complicado que otros factores me hagan probarlo.

	Las reseñas de otros jugadores y el desarrollador me son complementarios. Las reseñas me sirven para ver si el juego tiene problemas no evidentes a simple vista, como pueden ser problemas de rendimiento, el juego volviéndose tedioso en etapas más avanzadas, etc. Por otro lado, el desarrollador me importa si jugué juegos de ellos anteriormente, o si tiene muy mala fama.

	\item[\textbf{3. ¿Qué tanta importancia les das a las recomendaciones de videojuegos que recibís y qué hace que decidas seguirlas o ignorarlas?}]
	Le doy bastante importancia a las recomendaciones de videojuegos que recibo, especialmente si vienen de un amigo o persona cercana que conoce mis gustos.

	Siempre que me recomienden un videojuego voy a darle más tiempo que el que le daría si lo hubiera descubierto por mi cuenta. Como mínimo voy a ver el tráiler o un tramo de gameplay, y suelo jugarlo una hora o dos también.

	Si la recomendación viene de alguien que me conoce bien, puedo llegar a prestarle más tiempo todavía, continuando jugándolo incluso si el comienzo del videojuego no me atrae.

	\item[\textbf{4. Cuando una plataforma te recomienda un videojuego, ¿qué información debería mostrarte para que entiendas por qué fue elegido para vos?}]
	Considero que compararlo con otros videojuegos que haya jugado en el pasado sería la información más importante, pero los géneros del videojuego y el desarrollador también me parecen información notable.

	En mi opinión, los géneros pueden indicar similitudes más simples o fáciles de ver a simple vista, pero la estética y otros aspectos difíciles de definir en palabras de un videojuego me parecen complicados de clasificar mediante géneros.

	Por ende, veo la comparación con otros videojuegos de estéticas parecidas, o de desarrolladores que compartan dicha estética, como una manera de recomendarme un videojuego que se siente más personal.

	\item[\textbf{5. Si pudieras diseñar una herramienta ideal para ayudarte a elegir videojuegos, ¿qué tendría que hacer para que realmente la usaras?}]
	Si tuviera la capacidad de desarrollar una herramienta ideal, pondría el foco en la personalización de las recomendaciones según el usuario. No tiene sentido recomendarle un videojuego de puzzles a alguien que solo juega videojuegos FPS, ni tampoco recomendarle videojuegos AAA a un usuario que mayormente juega videojuegos indies. Diría que la falta de esta atención personal a cada usuario es lo que personalmente me hace buscar videojuegos por otros lados, fuera de las plataformas de distribución de videojuegos que existen hoy en día.
\end{description}
